\chapter{General Presentation}

\section{Introduction}

This chapter establishes the strategic and operational framing of the project. It presents the general framework, introduces the host organization, details project background and problematic, and studies the existing intervention management process to justify the proposed digital solution.

\section{General Framework of the Project}

The RM project belongs to the digital transformation of field operations in telecom and IT equipment installation contexts. In such environments, intervention quality depends on rapid coordination between office teams and field technicians, consistent data collection, and reliable validation cycles.

From a systems engineering perspective, this context imposes five structural constraints:

\begin{enumerate}
\item distributed operations across multiple sites,
\item heterogeneous intervention scenarios requiring configurable forms,
\item strict accountability expectations from clients and operators,
\item tight turnaround targets for validation and reporting,
\item and high dependence on mobile data capture under variable connectivity.
\end{enumerate}

Consequently, the project is not only a software implementation effort. It is also a process reengineering initiative where workflow governance, traceability, and data quality control are central design objectives.

\section{Presentation of the Host Organization}

\subsection{Organizational Identity}

Africom Services is a Tunisian company founded in 2005 and active in telecommunications engineering, field deployment, and maintenance services. The organization supports telecom operators, enterprise clients, and institutional actors through project execution and technical operations.

\subsection{Core Activity Areas}

\begin{itemize}
\item deployment and validation of telecom transmission equipment,
\item radio network and microwave infrastructure projects,
\item fiber optic implementation and maintenance,
\item preventive and corrective technical interventions,
\item operational monitoring solutions and connected security services.
\end{itemize}

\subsection{Operational Characteristics}

The organization combines office-based coordination with field execution teams. This dual operating model creates strong demand for synchronized information systems able to connect planning, execution evidence, quality control, and reporting.

\begin{figure}[H]
\centering
\fbox{\parbox{0.88\textwidth}{
\centering
% TODO: Insert host organization chart here
Organization Chart Placeholder: Operational management, coordination unit, supervision unit, field technician teams.
}}
\caption{Host Organization Structure Placeholder}
\label{fig:host_org_structure_placeholder}
\end{figure}

\section{Project Presentation}

\subsection{Project Background}

Before RM, intervention governance relied on partially fragmented practices with heterogeneous tools and significant manual effort. The operational flow from assignment to final report delivery lacked a unified digital backbone. As intervention volume increased, this model exposed scalability limits.

The PRD therefore defined RM as an integrated web and mobile platform with the following business mission:

\begin{itemize}
\item standardize intervention execution through template-driven processes,
\item improve supervisor visibility on field progress,
\item enforce validation and rejection governance,
\item and automate client-ready Excel reporting.
\end{itemize}

\subsection{Problematic}

The project problematic can be formalized as:

\begin{quote}
How can a telecom field-operations organization move from fragmented, partly manual intervention tracking to a controlled, auditable, and scalable digital workflow without reducing field agility?
\end{quote}

This problematic combines business and technical dimensions:

\begin{itemize}
\item business dimension: cycle time reduction, quality consistency, and managerial visibility,
\item technical dimension: secure multi-platform architecture, dynamic form support, and robust traceability.
\end{itemize}

\section{Study of the Existent}

\subsection{Analysis of the Existing}

The initial process can be summarized in six stages:

\begin{enumerate}
\item request preparation by coordination staff,
\item technician assignment through ad-hoc communication,
\item field execution and evidence collection,
\item supervisor quality checks,
\item correction loop when errors are detected,
\item final report consolidation for client communication.
\end{enumerate}

Observed process weaknesses were concentrated at handoff points between stakeholders, where information loss and delay frequently occurred.

\begin{table}[H]
\centering
\caption{Observed Weaknesses in the Existing Process}
\label{tab:existing_process_weaknesses}
\begin{tabularx}{\textwidth}{|l|X|X|}
\hline
\textbf{Process Stage} & \textbf{Observed Issue} & \textbf{Operational Impact} \\
\hline
Assignment and planning & Inconsistent data format between teams & Rework and interpretation errors \\
\hline
Field execution & Non-standardized capture structure & Data quality variability \\
\hline
Validation & Limited immediate access to complete evidence & Longer decision cycle time \\
\hline
Reporting & Manual consolidation effort & Delayed delivery and reduced scalability \\
\hline
Audit and compliance & Partial history tracking & Difficult incident analysis and accountability \\
\hline
\end{tabularx}
\end{table}

\subsection{Critique of the Existing}

The existing process can be critiqued using a quality-attribute lens:

\begin{itemize}
\item \textbf{Traceability deficit}: status transitions and decision rationale were not uniformly logged.
\item \textbf{Standardization deficit}: different technicians used different capture practices.
\item \textbf{Responsiveness deficit}: supervisors had delayed visibility on field outputs.
\item \textbf{Scalability deficit}: manual aggregation could not scale with intervention growth.
\item \textbf{Governance deficit}: control rules were implicit rather than encoded in system behavior.
\end{itemize}

Industry comparison confirms that modern field-service platforms prioritize event-driven status management, structured evidence collection, and automated report generation. The existing model did not provide these capabilities in an integrated way.

\subsection{Proposed Solution}

RM addresses the identified deficits through a domain-driven digital platform organized around five pillars:

\begin{enumerate}
\item \textbf{Identity and governance}: JWT authentication with role-based control.
\item \textbf{Workflow control}: lifecycle management from planning to approval and archival.
\item \textbf{Template-driven execution}: dynamic forms for standardized field data capture.
\item \textbf{Evidence and traceability}: files, geolocation, and status history.
\item \textbf{Automated reporting}: configurable Excel generation aligned with operator needs.
\end{enumerate}

\begin{figure}[H]
\centering
\fbox{\parbox{0.9\textwidth}{
\centering
% TODO: Insert as-is vs to-be process diagram here
Process Transformation Placeholder: Existing fragmented flow versus RM integrated digital workflow.
}}
\caption{As-Is and To-Be Process Comparison Placeholder}
\label{fig:as_is_to_be_placeholder}
\end{figure}

\section{Conclusion}

This chapter provided the business and organizational grounding of the RM project. It demonstrated why the legacy process was insufficient for current operational demands and justified the need for a structured digital platform. The next chapter presents the planning approach, tools, and architecture decisions used to operationalize this transformation.
